\documentclass[12pt]{article}
\usepackage[utf8]{inputenc}
\usepackage[T1]{fontenc}
\usepackage{hyperref}
\title{Chichicapan Zapotec Language Resources for K-12 Curriculum Development}
\author{Compiled by Claude Sonnet 5.5 (claude-sonnet-5-5) for Andr\'as Kornai}
\date{}
\begin{document}
\maketitle

\section{Introduction}
Chichicapan Zapotec (ISO 639-3 zpv; Glottocode chic1274) is an endangered Otomanguean language of San Baltazar Chichic\'apam, Oaxaca, Mexico. It is also called Eastern Ocotl\'an Zapotec. Wikipedia gives about 2,700 speakers (2005). The documentation found is archival, mostly elicitation recordings at AILLA, plus evangelistic audio. No school materials were found.

\section{Primary Educational Resources}
No primers or textbooks were found.

\section{Linguistic and Reference Materials}
\begin{itemize}
\item Wikipedia article on Chichic\'apam Zapotec: \url{https://en.wikipedia.org/wiki/Chichic%C3%A1pam_Zapotec}
\item AILLA, phonemes and examples in Chichicapan Central Zapoteco: \url{https://islandora-ailla.lib.utexas.edu/islandora/object/ailla%3A135293}
\item AILLA, illustrated dictionary elicitation in Chichicapan Central Zapoteco: \url{https://islandora-ailla.lib.utexas.edu/islandora/object/ailla%3A135249}
\item AILLA, minimal pairs: \url{https://islandora-ailla.lib.utexas.edu/islandora/object/ailla%3A135270}
\item AILLA, announcement of the Survey of Zapotec and Chatino Languages Collection (Sicoli and Kaufman): \url{https://islandora-ailla.lib.utexas.edu/node/151}
\end{itemize}
Access levels for some AILLA files may be restricted. I did not check them.

\section{Religious and Community Resources}
\begin{itemize}
\item AILLA, narrative: Changing of the Authorities in San Baltazar Chichicapan: \url{https://islandora-ailla.lib.utexas.edu/islandora/object/ailla%3A135263}
\end{itemize}

\section{Digital and Audio Resources}
\begin{itemize}
\item Global Recordings Network, audio: \url{https://www.globalrecordings.net/en/language/zpv}
\item Global Recordings Network, alternate page: \url{https://globalrecordings.net/en/language/4440}
\end{itemize}

\section{Curriculum Development Context}
The language is used as a first language by adults only, so intergenerational transmission is a concern. The AILLA recordings of lexical items, phonemes and narratives are the best starting material for a lexicon and an orthography. They would need community permission and access arrangements.

\end{document}